\documentclass[a4paper,12pt]{article}

\usepackage{ucs}
\usepackage[utf8x]{inputenc}
\usepackage[russian]{babel}
\usepackage{graphicx}
\usepackage{listings}
\usepackage{xcolor}
\usepackage{geometry}
\usepackage{indentfirst}
\usepackage{array}
\usepackage{float}
\usepackage{caption}
\usepackage{pdflscape}

\makeatletter
\renewcommand\@biblabel[1]{#1.}
\makeatother

\newcommand{\myrule}[1]{\rule{#1}{0.4pt}}
\newcommand{\sign}[2][~]{{\small\myrule{#2}\\[-0.7em]\makebox[#2]{\it #1}}}

% Поля
\geometry{top=20mm, left=30mm, right=10mm, bottom=20mm, nohead}

% Межстрочный интервал
\renewcommand{\baselinestretch}{1.50}

\begin{document}

%%%%%%%%%%%%%%%%%%%%%%%%%%%%%%%
%%%                         %%%
%%% Начало титульного листа %%%

\thispagestyle{empty}
\begin{center}
\renewcommand{\baselinestretch}{1}
{\large
{\sc Петрозаводский государственный университет\\
Институт математики и информационных технологий\\
	Кафедра ИМО
}}
\end{center}

\begin{center}
Направление подготовки бакалавриата \\
09.03.02 - Информационные системы и технологии
\end{center}

\vfill

\begin{center}
{\normalsize Отчет о практике по научно-исследовательской работе}

\medskip

%%% Название работы %%%
	{\Large \sc РАЗВИТИЕ СИСТЕМЫ "КУРС"} \\
	(промежуточный)
\end{center}

\medskip

\begin{flushright}
\parbox{11cm}{%
\renewcommand{\baselinestretch}{1.2}
\normalsize
	Выполнили:\\
студенты группы 22406 и 22405
\begin{flushright}
	Пянтин А. С. \sign[подпись]{4cm}
\end{flushright}
\begin{flushright}
	Титов Е. В. \sign[подпись]{4cm}
\end{flushright}

Место прохождения практики: \\ кафедра ИМО

Период прохождения практики: \\

Руководитель:\\
Доцент Кулаков Кирилл Александрович \\
\begin{flushright}
\sign[подпись]{4cm}
\end{flushright}

Итоговая оценка
\begin{flushright}
  \sign[оценка]{4cm}
\end{flushright}
}
\end{flushright}

\vfill

\begin{center}
\large
    Петрозаводск --- 2025
\end{center}

%%% Конец титульного листа  %%%
%%%                         %%%
%%%%%%%%%%%%%%%%%%%%%%%%%%%%%%%

\newpage
\tableofcontents

\newpage
\section*{Введение}
\addcontentsline{toc}{section}{Введение}

В настоящее время информационные технологии активно применяются в образовательной сфере для автоматизации различных процессов. Одним из важных направлений является разработка веб-приложений для управления учебным процессом, в том числе для организации работы с дипломными проектами студентов. Современные веб-приложения должны обеспечивать удобный интерфейс для всех участников образовательного процесса: студентов, преподавателей и администраторов. При этом важными требованиями являются интуитивность интерфейса, быстродействие системы, возможность работы без перезагрузки страниц, а также автоматизация рутинных операций.

В данной работе рассматривается модернизация веб-приложения для управления дипломными работами, разработанного на основе фреймворка Django. Приложение предназначено для организации процесса выбора тем дипломных работ студентами, их утверждения преподавателями, а также для управления связанными документами и данными. Целью работы является расширение функциональности веб-приложения для управления дипломными работами, улучшение пользовательского интерфейса и автоматизация административных задач.

Для достижения поставленной цели необходимо решить следующие задачи: расширить модель данных для тем дипломных работ, добавив новые поля и функциональность; модернизировать интерфейс страниц «Предложить тему» и «Выбрать тему» с использованием технологий AJAX и JavaScript для улучшения пользовательского опыта; реализовать систему уведомлений для преподавателей и студентов; разработать набор management-команд для автоматизации административных задач: переход на новый учебный год, генерация тестовых данных, создание групп; расширить возможности административной панели для управления данными системы.

В ходе выполнения работы были использованы следующие технологии и инструменты: Django — веб-фреймворк для разработки серверной части приложения, JavaScript (чистый) — для реализации интерактивности на клиентской стороне, AJAX — для асинхронного взаимодействия с сервером без перезагрузки страниц, CSS — для стилизации интерфейса, Faker — библиотека для генерации тестовых данных.

Практическая значимость работы определяется тем, что разработанные решения могут быть использованы в реальных образовательных учреждениях для автоматизации процесса управления дипломными работами, что позволит сократить время на выполнение рутинных операций и улучшить качество взаимодействия между студентами и преподавателями.

Данная работа состоит из введения, четырёх глав, заключения и списка литературы. В первой главе описывается общая архитектура приложения и существующая функциональность. Во второй главе рассматриваются изменения в моделях данных, формах и представлениях. В третьей главе описываются изменения в пользовательском интерфейсе и административных инструментах. В четвёртой главе представлены метрики выполненной работы, демонстрирующие равномерное распределение задач между участниками проекта. В заключении подводятся итоги проделанной работы, формулируются полученные результаты.

Работа выполнена двумя участниками: Пянтин Артем Сергеевич (группа 22406) — разработка моделей данных, форм, middleware, базовой части представлений для преподавателей и студентов, административной панели, management-команды roll\_year и части generate\_test\_data, а также шаблонов для управления темами преподавателя и части шаблонов для студентов. Титов Евгений Викторович (группа 22405) — разработка основной логики представлений для преподавателей и студентов, вспомогательных функций обработки заявок, большей части management-команды generate\_test\_data, management-команды create\_groups, URLs, а также основной части шаблонов для выбора тем студентами и связанных фрагментов.

\newpage
\section{Архитектура и существующая функциональность приложения}

\subsection{Общее описание системы}

Веб-приложение для управления дипломными работами представляет собой систему, разработанную на основе фреймворка Django. Приложение предназначено для организации процесса выбора и управления темами дипломных работ студентами и преподавателями. Основными участниками системы являются: студенты — могут просматривать доступные темы дипломных работ, подавать заявки на выбранные темы, просматривать статус своих заявок; преподаватели — могут создавать темы дипломных работ, просматривать заявки от студентов, утверждать или отклонять заявки, управлять своими темами; администраторы — имеют полный доступ к системе для управления пользователями, группами, годами обучения и другими данными.

Система использует модель данных, включающую следующие основные сущности: Student — информация о студентах, TeacherProfile — профили преподавателей, Year — учебные годы, Group — группы студентов, Enrollment — записи о зачислении студентов на учебные годы с информацией о дипломной работе, Topic — темы дипломных работ, TopicRequest — заявки студентов на темы, Document — документы, связанные с дипломными работами.

\newpage
\section{Структура базы данных системы}

\subsection{Общая архитектура данных}

Структура базы данных системы «КУРС» построена на основе реляционной модели с использованием Django ORM в качестве объектно-реляционного отображения. В процессе разработки применялась СУБД SQLite3, что обеспечивает простоту развертывания и отладки в условиях разработки. Архитектура данных спроектирована с учетом специфики образовательного процесса и включает 8 основных таблиц, а также стандартную таблицу пользователей Django \texttt{auth\_user}.

Система использует различные типы связей между таблицами: ForeignKey для организации отношений «многие-к-одному» и OneToOneField для создания уникальных связей между записями. Такая архитектура обеспечивает целостность данных и позволяет эффективно управлять информационным потоком между различными сущностями системы.

\subsection{Диаграмма связей между таблицами}
\begin{figure}[H]
\centering
\includegraphics[width=0.75\textwidth]{database_schema.png}
\caption{Диаграмма структуры базы данных системы}
\label{fig:db_schema}
\end{figure}

\subsection{Описание основных таблиц}

\subsubsection{Таблица Year (Учебный год)}

Таблица предназначена для хранения информации об учебных годах, что позволяет организовывать данные по временным периодам. Каждый учебный год представляет собой отдельный временной контекст, к которому привязываются группы студентов, записи обучения и связанные с ними данные. Поле \texttt{year} содержит целочисленное значение (например, 2024) и должно быть уникальным в пределах таблицы. Таблица Year служит основой для временной организации всех данных системы, обеспечивая возможность работы с несколькими учебными годами одновременно и хранения исторических данных.

\subsubsection{Таблица Group (Группа студентов)}

Данная таблица хранит информацию о группах студентов, привязанных к конкретным учебным годам. Каждая группа имеет уникальное название в рамках учебного года (например, «22305») и флаг \texttt{is\_latest}, указывающий на то, является ли группа выпускной (4-й курс бакалавриата). Ограничение \texttt{unique\_together} на поля \texttt{name} и \texttt{year} гарантирует, что в рамках одного учебного года не может существовать двух групп с одинаковым названием. Связь с таблицей Year реализована через ForeignKey с каскадным удалением, что означает автоматическое удаление всех групп при удалении соответствующего учебного года.

\subsubsection{Таблица Student (Студент)}

Таблица содержит базовую информацию о студентах системы. Поле \texttt{login} представляет собой уникальный идентификатор студента, а поле \texttt{full\_name} хранит полное имя. Такая структура позволяет однозначно идентифицировать каждого студента в системе. Связи с другими таблицами организованы через таблицу Enrollment, что обеспечивает нормализацию данных и предотвращает избыточность хранения информации.

\subsubsection{Таблица Enrollment (Запись обучения)}

Эта таблица является центральным связующим элементом системы, соединяющим студентов с учебными годами, группами и информацией о дипломных работах. Каждая запись представляет собой привязку конкретного студента к конкретному учебному году и группе. Поля, связанные с дипломной работой (такие как \texttt{title}, \texttt{adviser\_name}, \texttt{adviser\_position}, \texttt{department}), могут оставаться пустыми до момента выбора темы студентом. Ограничение \texttt{unique\_together} на поля \texttt{student} и \texttt{year} гарантирует, что каждый студент может иметь только одну запись обучения на учебный год.

\subsubsection{Таблица Document (Документ)}

Таблица предназначена для хранения загруженных документов, связанных с дипломными работами студентов. Система поддерживает различные типы документов, которые классифицируются по полю \texttt{doc\_type}. Для промежуточных документов доступны типы «Пр. отчет» и «Пр. ЭП». Для обычных финальных работ предусмотрены типы «Отчет» и «ЭП». Для выпускных групп (с флагом \texttt{is\_latest}) доступны дополнительные типы документов: «Отчет по практике НИР», «Текст ВКР», «Презентация ВКР», «Проверка на плагиат» и «Отзыв руководителя». Путь загрузки файлов формируется динамически на основе иерархии: \texttt{media/students/\{год\}/\{группа\}/\{студент\}/\{имя\_файла\}}.

\subsubsection{Таблица TeacherProfile (Профиль преподавателя)}

Таблица расширяет стандартную модель пользователя Django дополнительной информацией, специфичной для преподавателей. Связь с таблицей \texttt{auth\_user} реализована через OneToOneField, что обеспечивает уникальность соответствия между пользователем системы и профилем преподавателя. Поле \texttt{adviser\_position} содержит информацию о должности преподавателя с возможными значениями: «преподаватель», «ст. преподаватель», «доцент», «профессор», «зав. кафедрой», «другая». Эта информация автоматически копируется в запись Enrollment при утверждении заявки студента.

\subsubsection{Таблица Topic (Тема курсовой работы)}

В данной таблице хранятся темы дипломных работ, предлагаемые преподавателями. В рамках модернизации системы были добавлены новые поля: \texttt{course} (курс, для которого предназначена тема, значения от 1 до 6), \texttt{capacity} (максимальное количество студентов, которые могут работать над темой одновременно) и \texttt{direction} (направление подготовки). Поле \texttt{is\_active} позволяет временно деактивировать тему без её удаления из системы. Автоматически обновляемые поля \texttt{created\_at} и \texttt{updated\_at} отслеживают время создания и последнего изменения темы.

\subsubsection{Таблица TopicRequest (Заявка на тему)}

Таблица управляет процессом подачи и рассмотрения заявок студентов на темы дипломных работ. Каждая заявка связывает конкретную тему с конкретной записью обучения студента. Поле \texttt{status} отслеживает состояние заявки с возможными значениями: «ожидает решения», «принята», «отклонена». Поле \texttt{comment} позволяет преподавателю оставлять комментарии при принятии решения по заявке. Ограничение \texttt{unique\_together} на поля \texttt{topic} и \texttt{enrollment} предотвращает создание нескольких заявок одного студента на одну тему.

\subsection{Дополнительные таблицы Django}

\subsubsection{Таблица auth\_user (Пользователи системы)}

Стандартная таблица Django для аутентификации пользователей. Содержит основные поля для управления учетными записями: \texttt{username} (уникальное имя пользователя), \texttt{password} (хеш пароля), \texttt{email}, \texttt{first\_name}, \texttt{last\_name}, а также флаги доступа \texttt{is\_staff}, \texttt{is\_superuser}, \texttt{is\_active}. Таблица интегрирована с системой безопасности Django и обеспечивает базовую аутентификацию и авторизацию пользователей.

\subsection{Ключевые особенности структуры базы данных}

Архитектура базы данных системы «КУРС» обладает рядом важных особенностей, обеспечивающих её эффективность и надежность. Многоуровневая иерархия данных организована следующим образом: учебные годы содержат группы, которые в свою очередь содержат записи обучения студентов, к которым привязаны документы и заявки на темы. Параллельно существует иерархия пользователей: стандартные пользователи Django могут иметь связанные профили преподавателей, которые создают темы дипломных работ.

Система ограничений целостности данных реализована с учетом бизнес-логики приложения. Для зависимых данных (например, группы при удалении учебного года) используется каскадное удаление (CASCADE). Для критических данных (группы, к которым привязаны записи обучения) применяется защита от удаления (PROTECT), что предотвращает случайную потерю важной информации. Ограничения \texttt{unique\_together} на различных уровнях обеспечивают уникальность комбинаций полей, предотвращая дублирование записей.

Гибкость структуры данных проявляется в нескольких аспектах. Поля таблицы Enrollment, связанные с дипломной работой, могут оставаться пустыми до момента выбора темы студентом, что соответствует реальному учебному процессу. Система поддерживает различные типы документов в зависимости от типа группы (выпускная или невыпускная). Механизм статусов заявок позволяет гибко управлять процессом утверждения тем.

Масштабируемость архитектуры обеспечивается поддержкой работы с несколькими учебными годами одновременно, возможностью хранения исторических данных без потери производительности и эффективной группировкой данных по кафедрам и направлениям подготовки. Такая структура позволяет системе адаптироваться к изменяющимся требованиям и увеличивающемуся объему данных.

Вопросы безопасности решены на нескольких уровнях. Связь преподавателей со стандартной системой аутентификации Django обеспечивает надежную защиту учетных записей. Уникальные логины студентов предотвращают конфликты идентификаторов. Использование ограничения PROTECT для критических таблиц защищает от случайного удаления важных данных.

\subsection{Статистические характеристики структуры}

Общая структура базы данных включает 8 основных таблиц с 34 основными полями. Между таблицами установлено 9 связей ForeignKey и OneToOneField. Для оптимизации производительности создано 10 индексов, включая составные индексы для часто используемых комбинаций полей. Система уникальных ограничений включает 7 различных комбинаций, обеспечивающих целостность данных на уровне базы данных.

Такая структура обеспечивает баланс между нормализацией данных для минимизации избыточности и денормализацией для оптимизации производительности часто выполняемых запросов. Архитектура спроектирована с учетом типичных сценариев использования системы и обеспечивает эффективную работу как при малых, так и при больших объемах данных.

\newpage
\section{Модернизация пользовательского интерфейса и административных инструментов}
\subsection{Сравнительная таблица изменений интерфейса и функциональности}

\begin{table}[H]
\centering
\caption{Сравнение функциональности старой и новой версии системы (часть 1)}
\begin{small}
\begin{tabular}{|p{0.22\textwidth}|p{0.33\textwidth}|p{0.33\textwidth}|}
\hline
\textbf{Аспект системы} & \textbf{Старая версия} & \textbf{Новая версия} \\
\hline
\hline
Модель данных темы & Только основные поля: название, описание, кафедра & Добавлена информация о курсе, количестве доступных мест и направлении подготовки \\
\hline
Профиль преподавателя & Только основная информация о преподавателе & Добавлено поле для указания должности преподавателя \\
\hline
Управление темами преподавателем & Форма с перезагрузкой страницы после каждого действия & Интерактивный интерфейс с возможностью массового добавления, редактирования и удаления тем без перезагрузки \\
\hline
Группировка тем & Отсутствовала группировка, все темы отображались списком & Интеллектуальная группировка: для преподавателей — по доступности и курсам, для студентов — по кафедрам и направлениям \\
\hline
Система уведомлений & Отсутствовали какие-либо уведомления & Визуальные счетчики: для преподавателей — количество ожидающих заявок, для студентов — количество доступных тем \\
\hline
Выбор тем студентами & Простой список с необходимостью перезагрузки страницы & Умный интерфейс с фильтрацией по курсу, защитой от повторных заявок и анимированными элементами \\
\hline
\end{tabular}
\end{small}
\label{tab:interface_comparison_1}
\end{table}

\begin{table}[H]
\centering
\caption{Сравнение функциональности старой и новой версии системы (часть 2)}
\begin{small}
\begin{tabular}{|p{0.22\textwidth}|p{0.33\textwidth}|p{0.33\textwidth}|}
\hline
\textbf{Аспект системы} & \textbf{Старая версия} & \textbf{Новая версия} \\
\hline
\hline
Автоматизация административных задач & Все задачи выполнялись вручную через административную панель & Добавлены команды для автоматического перехода на новый учебный год, создания групп и генерации тестовых данных \\
\hline
Административная панель & Базовая функциональность управления данными & Расширенные возможности поиска, фильтрации и отображения данных \\
\hline
Клиентские технологии & Статические страницы с перезагрузками & Динамические страницы с асинхронными запросами, анимациями и современным дизайном элементов \\
\hline
Инструменты разработки & Базовый набор инструментов Django & Добавлен инструмент для генерации реалистичных тестовых данных \\
\hline
Обработка заявок студентов & Базовая логика утверждения и отклонения заявок & Автоматическое заполнение данных о преподавателе, интеллектуальное управление связями между сущностями \\
\hline
\end{tabular}
\end{small}
\label{tab:interface_comparison_2}
\end{table}
\subsection{Детальное описание улучшений интерфейса и функциональности}

Модернизация интерфейса страницы управления темами преподавателя (page\_teacher\_topics.html) была проведена с использованием современных веб-технологий. Была реализована расширяемая таблица для ввода новых тем с кнопкой «+» для добавления новых строк ввода, возможностью ввода нескольких тем одновременно через действие create\_batch, сохранением тем без перезагрузки страницы через AJAX. Существующие темы отображаются в режиме просмотра с возможностью редактирования: иконки редактирования/удаления для каждой темы, inline-редактирование полей темы прямо в таблице через действие update\_topic, AJAX-сохранение изменений. Реализована группировка тем для удобной навигации: группа «Свободные темы» отображает темы с доступными местами для студентов, группировка по курсам — темы сгруппированы по номерам курсов (1-6). Важным улучшением стала обработка удаления тем: при вызове действия delete\_topic происходит автоматическая отвязка всех студентов, работающих над темой, и очистка связанных данных в модели Enrollment.

Модернизация интерфейса выбора тем студентами (page\_topic\_select.html) и связанного фрагмента (topic\_table\_fragment.html) была значительно улучшена. Темы автоматически фильтруются по курсу студента, определяемому из его записи Enrollment для текущего года. При выборе всех кафедр темы группируются по кафедре и по направлению подготовки. Реализована проверка, запрещающая студенту подавать несколько ожидающих заявок одновременно. Кнопка подачи заявки была стилизована в стиле Bootstrap без подключения всего фреймворка. Реализована анимация раскрытия описания темы с использованием JavaScript и requestAnimationFrame: плавное раскрытие/сворачивание описания при клике на тему. Добавлена визуальная индикация доступных мест: отображение количества свободных и занятых мест, блокировка кнопки «Подать заявку» при достижении лимита утверждённых студентов.

Создан переиспользуемый фрагмент строки темы (topic\_row\_fragment.html) для страницы преподавателя, что обеспечивает единообразие интерфейса и упрощает поддержку кода. В шаблоне user\_header.html были внесены важные изменения: увеличены ограничения на длину логина и пароля для поддержки более длинных значений, добавлены счётчики уведомлений рядом с пунктами меню — счётчик ожидающих заявок для преподавателей рядом с «Предложить тему» и счётчик доступных тем для студентов рядом с «Выбрать тему».

Внесены улучшения в следующие шаблоны: page\_identity.html — улучшено форматирование и отображение информации о студентах и их дипломных работах; page\_query.html — оптимизированы запросы и улучшено отображение результатов поиска; report\_template.html — улучшено форматирование отчётов, добавлена поддержка новых полей (курс, направление, должность преподавателя).

Была добавлена новая зависимость в файл requirements.txt: Faker>=19.0.0 — для генерации тестовых данных, что значительно упростило процесс разработки и тестирования.

После модернизации система поддерживает следующие улучшенные сценарии. Сценарий работы преподавателя: пакетное добавление тем без перезагрузки страницы через create\_batch, inline-редактирование тем с мгновенным сохранением через update\_topic, удаление тем с автоматической отвязкой студентов через delete\_topic, группировка тем для удобной навигации (свободные и по курсам), получение уведомлений о новых заявках через middleware, удобное управление заявками студентов через действие decision. Сценарий работы студента: просмотр тем, отфильтрованных по своему курсу, группировка тем по кафедрам и направлениям при выборе всех кафедр, видимость доступных мест и занятых студентов, плавное раскрытие описаний тем с анимациями, защита от подачи множественных заявок, стилизованная кнопка подачи заявки, уведомления о новых доступных темах через middleware. Сценарий работы администратора: автоматическое создание групп через команду create\_groups, генерация тестовых данных для разработки и тестирования через generate\_test\_data, автоматизированный переход на новый учебный год с очисткой сессий через roll\_year, расширенные возможности админки с фильтрами и поиском.

Ключевым техническим улучшением стало внедрение AJAX-взаимодействия через JavaScript fetch API, что позволило устранить необходимость перезагрузки страниц при большинстве действий. Использование requestAnimationFrame для анимаций обеспечило плавность интерфейса. Middleware FoundYearMiddleware реализовал интеллектуальную систему уведомлений, вычисляющую релевантные данные для каждого типа пользователя. Management-команды автоматизировали наиболее трудоемкие административные задачи, такие как переход на новый учебный год и создание тестовых данных.

Все описанные улучшения значительно повысили удобство использования системы, автоматизировали рутинные административные задачи и создали современный пользовательский интерфейс, соответствующий актуальным веб-стандартам.

\newpage
\section{Метрики выполненной работы}

В данной главе представлены метрики выполненной работы, демонстрирующие равномерное распределение задач между двумя участниками проекта. Работа была разделена таким образом, чтобы каждый участник работал как над серверной, так и над клиентской частью приложения, с максимально равномерным распределением объёма работы.

\subsection{Распределение работы между участниками}

Работа была разделена между двумя участниками следующим образом: Пянтин Артем Сергеевич (группа 22406) — разработка моделей данных, форм, middleware, части представлений для преподавателей, части представлений для студентов, административной панели, management-команды roll\_year и части generate\_test\_data, а также шаблонов для управления темами преподавателя и части шаблонов для студентов. Титов Евгений Викторович (группа 22405) — разработка основной части представлений для преподавателей, основной части представлений для студентов, части management-команды generate\_test\_data, management-команды create\_groups, URLs, а также основной части шаблонов для выбора тем студентами и связанных фрагментов.

\subsection{Метрики работы участников}

\begin{table}[H]
\centering
\caption{Метрики выполненной работы по участникам}
\begin{tabular}{|l|c|c|}
\hline
\textbf{Параметр} & \textbf{Пянтин А.С.} & \textbf{Титов Е.В.} \\
\hline
Всего строк кода (чистое изменение) & 1251 & 1278 \\
\hline
Backend строк & 541 & 749 \\
\hline
Frontend строк & 710 & 529 \\
\hline
Количество измененных/созданных файлов & 7 & 6 \\
\hline
Количество функций/методов & 10 & 7 \\
\hline
Management-команды & 2 (1 полная + 1/2) & 2 (1 полная + 1/2) \\
\hline
HTML-шаблоны & 2 & 4 \\
\hline
JavaScript функции & 3 & 2 \\
\hline
CSS строк & ~210 & ~185 \\
\hline
\end{tabular}
\label{tab:participant_metrics}
\end{table}

\subsection{Детальное распределение по типам задач}

\begin{table}[H]
\centering
\caption{Детальное распределение задач по участникам (строки чистого изменения)}
\begin{tabular}{|l|c|c|}
\hline
\textbf{Задача} & \textbf{Пянтин А.С.} & \textbf{Титов Е.В.} \\
\hline
Модели данных (models.py) & +101 & --- \\
\hline
Формы (forms.py) & +10 & --- \\
\hline
Middleware (middleware.py) & +36 & --- \\
\hline
Views (преподаватели, базовая часть) & +110 & --- \\
\hline
Views (преподаватели, основная логика) & --- & +170 \\
\hline
Views (студенты, базовая часть) & +110 & --- \\
\hline
Views (студенты, основная часть) & --- & +226 \\
\hline
Вспомогательные функции views & --- & +40 \\
\hline
Административная панель (admin.py) & +24 & --- \\
\hline
Команда roll\_year & +129 & --- \\
\hline
Команда generate\_test\_data (часть) & +130 & +218 \\
\hline
Команда create\_groups & --- & +93 \\
\hline
URLs (urls.py) & --- & +2 \\
\hline
Зависимости (requirements.txt) & +1 & --- \\
\hline
Шаблон page\_teacher\_topics.html & +591 & --- \\
\hline
Фрагмент topic\_row\_fragment.html & +119 & --- \\
\hline
Шаблон page\_topic\_select.html & --- & +376 \\
\hline
Фрагмент topic\_table\_fragment.html & --- & +100 \\
\hline
Шаблон user\_header.html & --- & +8 \\
\hline
Шаблон report\_template.html & --- & +36 \\
\hline
Шаблон page\_query.html & --- & +9 \\
\hline
Шаблон page\_identity.html & --- & 0 \\
\hline
\end{tabular}
\label{tab:detailed_distribution}
\end{table}

\subsection{Анализ равномерности распределения}

Анализ метрик показывает следующее распределение работы. Общий объём Backend изменений составляет 1290 строк, из которых Пянтин А.С. выполнил 541 строку (41.9\%), а Титов Е.В. — 749 строк (58.1\%). Распределение Backend равномерное, при этом Титов Е.В. выполнил основную логику представлений для преподавателей и студентов, а также большую часть генерации тестовых данных, что объясняет его больший объём работы в этой части. Общий объём Frontend изменений составляет 1239 строк, из которых Пянтин А.С. выполнил 710 строк (57.3\%), а Титов Е.В. — 529 строк (42.7\%). Распределение Frontend отражает логическое разделение: Пянтин А.С. работал над более сложным интерфейсом управления темами преподавателя с AJAX-функциональностью, а Титов Е.В. — над интерфейсом выбора тем студентами с анимациями. Общий объём изменений составляет 2529 строк, из которых Пянтин А.С. выполнил 1251 строку (49.5\%), а Титов Е.В. — 1278 строк (50.5\%). Разница составляет всего 27 строк, что свидетельствует о практически равном вкладе участников.

\subsection{Выводы по метрикам}

Анализ метрик показывает, что оба участника внесли существенный и практически равный вклад в проект. Пянтин А.С. выполнил разработку Backend (541 строка): модели данных (+101), формы (+10), middleware (+36), базовая часть представлений для преподавателей и студентов (+110), административная панель (+24), management-команда roll\_year (+129), часть management-команды generate\_test\_data (+130), зависимости (+1). Frontend (710 строк): шаблоны для управления темами преподавателя (page\_teacher\_topics.html +591, topic\_row\_fragment.html +119). Всего: 1251 строка чистого изменения.

Титов Е.В. выполнил разработку Backend (749 строк): основная логика представлений для преподавателей (+170), основная часть представлений для студентов (+226), вспомогательные функции обработки заявок (+40), management-команда create\_groups (+93), большая часть management-команды generate\_test\_data (+218), URLs (+2). Frontend (529 строк): полный шаблон выбора тем студентами (page\_topic\_select.html +376), фрагмент таблицы тем (topic\_table\_fragment.html +100), заголовок пользователя (user\_header.html +8), другие шаблоны (report\_template.html +36, page\_query.html +9). Всего: 1278 строк чистого изменения.

Оба участника работали как над серверной, так и над клиентской частью приложения, что обеспечило комплексное понимание системы каждым из них. Распределение работы отражает логическое разделение функциональности с перекрытием: оба участника работали над представлениями для преподавателей и студентов, а также над шаблоном выбора тем, что обеспечило согласованность интерфейса и бизнес-логики. В Backend части Титов Е.В. выполнил больше работы (749 строк против 541), что объясняется реализацией основной бизнес-логики представлений и большей части генерации тестовых данных. В Frontend части Пянтин А.С. выполнил больше работы (710 строк против 529), что объясняется реализацией сложного интерфейса управления темами с AJAX-функциональностью. Оба участника внесли существенный и практически равный вклад в проект, работая над взаимодополняющими частями системы, что обеспечило целостность и функциональность разработанного решения.

\newpage
\section*{Заключение}
\addcontentsline{toc}{section}{Заключение}

В результате выполненной работы было успешно проведено развитие системы «КУРС» — веб-приложения для управления дипломными работами на основе фреймворка Django. Работа выполнялась двумя участниками: Пянтиным Артемом Сергеевичем (группа 22406) и Титовым Евгением Викторовичем (группа 22405), которые внесли практически равный вклад в разработку как серверной, так и клиентской части приложения. Основные достижения работы включают расширение модели данных для тем дипломных работ с добавлением полей course, capacity и direction, что позволило улучшить процесс фильтрации и группировки тем, а также контролировать количество студентов, работающих над одной темой. Расширение модели TeacherProfile полем adviser\_position автоматизировало сохранение должности преподавателя при утверждении заявки студента, упростив формирование отчетов и документов. Реализация системы уведомлений через middleware FoundYearMiddleware обеспечила подсчет ожидающих заявок для преподавателей и доступных тем для студентов с отображением счетчиков в интерфейсе.

Значительные улучшения были внесены в пользовательский интерфейс. Страница «Предложить тему» для преподавателей была переработана с использованием AJAX-технологий, что позволило реализовать пакетное добавление тем, inline-редактирование, удаление тем с автоматической отвязкой студентов и группировку тем по различным критериям. Страница «Выбрать тему» для студентов была улучшена за счет внедрения группировки тем по кафедрам и направлениям, фильтрации по курсу студента, плавных анимаций раскрытия описаний, визуальной индикации доступных мест и защиты от подачи множественных заявок. Интерфейсные элементы были стилизованы в современном дизайне с использованием CSS-анимаций и JavaScript-эффектов.

Для автоматизации административных задач были разработаны три management-команды. Команда roll\_year автоматизирует переход на новый учебный год с переносом групп, очисткой данных и сбросом сессий. Команда generate\_test\_data генерирует реалистичные тестовые данные с использованием библиотеки Faker для упрощения процесса разработки и тестирования. Команда create\_groups создает группы студентов с настраиваемыми параметрами, что значительно экономит время администраторов при подготовке учебного года.

Административная панель Django Admin была расширена с настройкой list\_display, фильтров и поиска для моделей TeacherProfile, Topic и TopicRequest, что упростило управление данными системы. Все внесенные изменения обеспечили обратную совместимость с существующими данными и функциональностью, что позволило внедрить улучшения без нарушения работы приложения.

Сравнение с предыдущей версией системы показывает значительный прогресс в следующих аспектах: пользовательский опыт улучшен за счет внедрения AJAX-технологий, что устранило необходимость перезагрузки страниц при большинстве действий; функциональность расширена новыми возможностями, такими как группировка тем, фильтрация по курсу, контроль количества мест; автоматизация административных задач сократила время на выполнение рутинных операций; интерфейс стал более современным и интуитивно понятным благодаря использованию анимаций, улучшенной обратной связи и визуальных индикаторов.

Практическая значимость работы заключается в том, что разработанные решения могут быть использованы в реальных образовательных учреждениях для автоматизации процесса управления дипломными работами. Внедрение описанных улучшений позволит сократить время на выполнение рутинных операций администраторами, улучшить качество взаимодействия между студентами и преподавателями, повысить удобство использования системы за счет современного интерфейса, автоматизировать процессы перехода на новый учебный год и упростить процесс тестирования и разработки за счет генерации тестовых данных.

В дальнейшем планируется развитие системы в следующих направлениях: добавление системы email-уведомлений для более оперативного информирования пользователей, расширение функциональности экспорта данных в различные форматы, реализация мобильной версии интерфейса для работы с мобильных устройств, добавление аналитики и отчетов по использованию системы, интеграция с другими образовательными системами.

Все поставленные задачи выполнены в полном объеме, система успешно модернизирована и готова к использованию в образовательном процессе. Распределение работы между участниками было сбалансированным: Пянтин А.С. внес 1251 строку чистого изменения (49.5\% от общего объема), Титов Е.В. — 1278 строк (50.5\%). Такой подход обеспечил комплексную проработку всех аспектов системы и достижение поставленных целей проекта.

\newpage
\renewcommand{\refname}{Библиографический список использованной литературы}
\begin{thebibliography}{99}
\vspace{5mm}
\addcontentsline{toc}{section}{Библиографический список использованной литературы}

\bibitem{django-docs}
Django Documentation [Электронный ресурс]. Режим доступа:
https://docs.djangoproject.com/

\bibitem{javascript-info}
JavaScript.info --- современный учебник JavaScript [Электронный ресурс]. Режим доступа:
https://javascript.info/

\bibitem{faker-docs}
Faker Documentation [Электронный ресурс]. Режим доступа:
https://faker.readthedocs.io/

\bibitem{ajax-guide}
MDN Web Docs: AJAX [Электронный ресурс]. Режим доступа:
https://developer.mozilla.org/en-US/docs/Web/Guide/AJAX

\bibitem{css-animations}
MDN Web Docs: CSS Animations [Электронный ресурс]. Режим доступа:
https://developer.mozilla.org/en-US/docs/Web/CSS/CSS\_Animations

\bibitem{django-admin}
Django Admin Documentation [Электронный ресурс]. Режим доступа:
https://docs.djangoproject.com/en/stable/ref/contrib/admin/

\bibitem{django-middleware}
Django Middleware Documentation [Электронный ресурс]. Режим доступа:
https://docs.djangoproject.com/en/stable/topics/http/middleware/

\bibitem{django-management-commands}
Django Management Commands [Электронный ресурс]. Режим доступа:
https://docs.djangoproject.com/en/stable/howto/custom-management-commands/

\bibitem{html5-details}
HTML5 details element [Электронный ресурс]. Режим доступа:
https://developer.mozilla.org/en-US/docs/Web/HTML/Element/details

\bibitem{request-animation-frame}
MDN Web Docs: requestAnimationFrame [Электронный ресурс]. Режим доступа:
https://developer.mozilla.org/en-US/docs/Web/API/window/requestAnimationFrame

\end{thebibliography}

\end{document}